\section{The global multiplication failure locus}\label{sec:global-geometry}
Throughout this section the field is $\C$, $k=n+1$, $4\le c<k$, $q=\binom{c+1}{2}$, and $p=2k-1$. On $X_L=\Gr(c,V_n)^L$ let $\mathcal S_\nu$ denote the tautological bundles. The failure scheme is defined by the maximal minors of
\[
 \mu:\bigoplus_{\nu=1}^L\Sym^2\mathcal S_\nu
                       \longrightarrow V_{2n}\otimes\mathcal O_{X_L}.
\]
Its underlying set is $\Delta^{(L)}$. We shall prove Theorem~\ref{thm:global-intro} without using a maximal-rank theorem for general curves.

\subsection{Annihilating Hankel forms}
For $\ell\in V_{2n}^*$, let $H_\ell(f,g)=\ell(fg)$ on $V_n$. The following classical input concerns the whole rank locus, including confluent and limiting forms; it is not an assertion only about sums of distinct point evaluations.
\begin{lemma}[Classical Hankel rank strata]\label{lem:hankel-strata}
In $\Pj(V_{2n}^*)$, the rank-exactly-$r$ stratum of $H_\ell$ is nonempty and has dimension
\[
 h_r=\begin{cases}2r-1,&1\le r<k,\\2k-2,&r=k.\end{cases}
\]
For $r<k$, its closure is the variety of sums of at most $r$ point evaluations, with Zariski closure understood.
\end{lemma}
\begin{proof}[Identification with the classical statement]
The rank-$\le r$ equations are the $(r+1)$-minors of the $k\times k$ Hankel matrix in $2k-1$ independent coordinates. The equality of Hankel minor ideals under changing the rectangular shape, recalled in \cite[Section 1]{CMSV}, identifies their coordinate ring with that of the maximal minors of an $(r+1)\times(2k-r-1)$ Hankel matrix. That ring has affine dimension $2r$. The secant identification is \cite[Section 2, equation (2.0.1)]{CMSV}; thus the projective dimension is $2r-1$. The rank-$\le r-1$ subvariety has smaller dimension, so the exact-rank stratum is nonempty and dense. The full-rank stratum is open in $\Pj^{2k-2}$. These are classical Hankel determinantal facts; no claim about their rational singularities is needed here.
\end{proof}

\begin{lemma}[Isotropic incidence dimensions]\label{lem:isotropic}
Let a symmetric form on a $k$-space have rank $r$ and radical $E$ of dimension $k-r$. For a totally isotropic $c$-plane $U$, put
\[
 s=\dim((U+E)/E),\qquad a=c-s.
\]
Such planes exist precisely for
\begin{equation}\label{eq:feasible}
 \max\{0,c-k+r\}\le s\le\min\{c,\lfloor r/2\rfloor\}.
\end{equation}
The stratum with this value of $s$ has codimension
\begin{equation}\label{eq:D-rs}
 D_{r,s}=(c-s)(r-s)+\frac{s(s+1)}2
\end{equation}
in $\Gr(c,k)$.
\end{lemma}
\begin{proof}
Choose $U\cap E$, an $a$-plane in $E$, with dimension $a(k-r-a)$. The image of $U$ in the nondegenerate $r$-dimensional quotient is an isotropic $s$-plane, whose Grassmannian has dimension $s(r-s)-s(s+1)/2$. This follows by differentiating the symmetric restriction equations at an isotropic frame: pairing with a complementary dual frame gives all symmetric variations, so the equations are independent. Existence is equivalent to $2s\le r$. Finally the lifts are an affine space of dimension $s(k-r-a)$. Subtracting the sum from $c(k-c)$ gives \eqref{eq:D-rs}. Maximal orthogonal Grassmannians may have two components; both have this same dimension, which is all the count uses.
\end{proof}

\begin{lemma}[Rank optimization]\label{lem:optimization}
Assume $Lq\ge2k-1$. For all $r,s$ satisfying \eqref{eq:feasible},
\begin{equation}\label{eq:optimization}
 L D_{r,s}-h_r\ge\min\{Lq-2k+2,Lc-3\}.
\end{equation}
\end{lemma}
\begin{proof}
Write $A_0=Lq-2k+2$ and $B_0=Lc-3$. If $r=k$, feasibility forces $s=c$, giving exactly $A_0$. Suppose $r<k$. For $s=0$, the left side is $(Lc-2)r+1\ge Lc-1>B_0$.

For $1\le s\le c-2$, the coefficient of $r$ in $L D_{r,s}-2r+1$ is $L(c-s)-2\ge0$. Since $r\ge2s$, the expression is at least
\[
 g(s)=1+s(Lc-4)-\frac{Ls(s-1)}2.
\]
This is concave as a function of $s$, and
\[
 g(1)=B_0,\qquad
 g(c-2)-B_0=(c-3)\left(\frac{L(c+2)}2-4\right).
\]
The latter is nonnegative except when $(L,c)=(1,4)$ or $(1,5)$. In the first exception, saturation forces $k=5$ and feasibility forces $s\le1$. In the second, $k\le8$; the values $g(1),g(2),g(3)$ are $2,2,1$. The last value is feasible only when $k=8$, in which case $A_0=1$. Both exceptions therefore satisfy \eqref{eq:optimization}.

For $s=c-1$ and $L=1$, the expression decreases with $r$, hence is at least its value at $r=k-1$, namely $A_0+k-2c+1$. Feasibility implies $k\ge2c-1$. For $s=c-1$ and $L\ge2$, it is minimized at $r=2c-2$ and exceeds $B_0$ by
\[
 (c-2)\left(\frac{L(c+1)}2-4\right)\ge0.
\]
Finally, if $s=c$ and $r<k$, it is $Lq-2r+1\ge A_0+1$. These cases exhaust \eqref{eq:feasible}.
\end{proof}

\subsection{The common-secant family}
Let $C_n\subset\Pj(V_n^*)$ be the evaluation rational normal curve, and let $S_n$ be the closure of its secant lines. Since $n\ge4$, $S_n$ is irreducible of dimension three. Define
\begin{equation}\label{eq:secant-incidence}
 \mathfrak S_L=\{([h],U_1,\ldots,U_L):[h]\in S_n,
                                  \ U_\nu\subset\ker h\},
 \qquad \Sigma_L=\operatorname{pr}_{X_L}\mathfrak S_L.
\end{equation}
\begin{proposition}[A global secant obstruction]\label{prop:secant}
The closed set $\Sigma_L$ is irreducible, has codimension $Lc-3$ in $X_L$, and lies in $\Delta^{(L)}$.
\end{proposition}
\begin{proof}
The incidence is a relative product of Grassmannians over $S_n$, so it is irreducible of dimension $3+Lc(k-1-c)$ and its projection is closed. For a general $h=\alpha\operatorname{ev}_x+\beta\operatorname{ev}_y$, with $x\ne y$ and $\alpha\beta\ne0$, every $f\in\ker h$ satisfies $\alpha f(x)=-\beta f(y)$. Consequently
\[
 \ell=\alpha^2\operatorname{ev}_x^{(2n)}-
                     \beta^2\operatorname{ev}_y^{(2n)}
\]
is nonzero and annihilates all products from each $U_\nu$. Closedness extends this conclusion to the whole incidence.

It remains to check the dimension of the image, rather than merely that of the incidence. Fix a general $h$. If $c=k-1$, the hyperplane $U_1$ determines $h$. Otherwise, the $c$-planes in $\ker h$ that are also annihilated by some $h'\in S_n\setminus\{h\}$ form a constructible set of dimension at most
\[
 3+c(k-2-c)<c(k-1-c),
\]
because $c\ge4$. Here, for each $h'$, the plane lies in the intersection of two distinct hyperplanes. Thus a general $U_1\subset\ker h$ admits no such $h'$, and the incidence projection is generically one-to-one. Subtraction from $\dim X_L=Lc(k-c)$ gives $Lc-3$.
\end{proof}

\begin{proof}[Proof of Theorem~\ref{thm:global-intro}]
Failure is equivalent to the existence of a nonzero $\ell$ annihilating every $U_\nu^2$. Consider the closed projective incidence of such pairs $((U_\nu),[\ell])$. On the rank-$r$ stratum of $H_\ell$, stratify each $U_\nu$ by $s_\nu$. Lemmas~\ref{lem:hankel-strata} and \ref{lem:isotropic} bound the incidence dimension by
\[
 \dim X_L+h_r-\sum_\nu D_{r,s_\nu}.
\]
Its image cannot have larger dimension. For fixed $r$, each term is at least the smallest feasible $D_{r,s}$, so Lemma~\ref{lem:optimization} gives
\[
 \codim\Delta^{(L)}\ge\min\{Lq-p+1,Lc-3\}\ge1.
\]
This proves, in particular, generic surjectivity without assuming it.

Proposition~\ref{prop:secant} gives the upper bound $Lc-3$ and shows that the failure set is nonempty. On a Grassmannian chart the map $\mu$ is a $p\times Lq$ matrix over a regular ring. The classical maximal-minor height bound gives codimension at most $Lq-p+1$ for its nonempty failure locus; see \cite[Introduction]{EHU}, with generic rank $p$ and minor size $p$. The two upper bounds match the lower bound. When $Lc-3$ is the minimum, the irreducible closed subvariety $\Sigma_L$ has the maximal dimension of $\Delta^{(L)}$, so it is an irreducible component of its underlying reduced set. If $Lq<p$, the dimension of the domain of $\mu$ makes surjectivity impossible. This proves the final equivalence.
\end{proof}

\begin{corollary}[The base-point locus and the contact count]\label{cor:basepoint}
For $L=1$, the locus $\mathcal B$ of spaces with a common projective zero is irreducible of codimension $c-1$, has multiplication corank at least two, and is properly contained in $\Sigma_1$. It is not an irreducible component of $\Delta^{(1)}$. For arbitrary $L$, the least number of independently chosen $c$-planes with full product span is
\[
 L_{\min}=\left\lceil\frac{2k-1}{\binom{c+1}{2}}\right\rceil.
\]
\end{corollary}
\begin{proof}
For a fixed linear factor $l$, common-base-point spaces form $\Gr(c,lV_{n-1})$. Varying $[l]\in\Pj^1$ gives an irreducible incidence of dimension $1+c(k-1-c)$, generically one-to-one because a general residual space has no further common factor. The codimension is $c-1$. Products are divisible by $l^2$, whose multiples have dimension $p-2$. The annihilating hyperplane is an evaluation point of $C_n\subset S_n$, proving inclusion in $\Sigma_1$. The codimensions $c-1$ and $c-3$ make that inclusion proper. An irreducible proper subset of the irreducible closed set $\Sigma_1\subset\Delta^{(1)}$ cannot itself be a component of $\Delta^{(1)}$. The contact count follows from Theorem~\ref{thm:global-intro} and the dimension obstruction.
\end{proof}
The theorem determines the exact dimension of the entire bad set and one explicit maximal-dimensional component in the stated regime. It does not assert that every component is a secant component, nor that the determinantal scheme is reduced, equidimensional, or everywhere transverse. In particular, the equality case of the two codimension bounds does not exclude additional components of the same dimension.
